\section{Encoder-Decoder 预训练：T5 与片段腐蚀}

\subsection{Encoder-Decoder 的预训练目标}

Encoder-Decoder 架构可以简单地用\textbf{语言建模}来预训练：将文本前半部分交给 Encoder（提供双向上下文），用 Decoder 自回归生成后半部分。

\begin{figure}[H]
\centering
\includegraphics[width=0.95\textwidth]{slides-images/slide-040.jpg}
\caption{Encoder-Decoder 的语言建模预训练：Encoder 处理前缀 $w_1, \ldots, w_T$（双向上下文），Decoder 自回归生成 $w_{T+1}, \ldots, w_{T+S}$\protect\footnotemark}
\end{figure}
\footnotetext{Slides 第41页。来自 Raffel et al., 2018。}

但在实践中，一种更有效的方法是\textbf{片段腐蚀}（Span Corruption）。

\subsection{T5 与片段腐蚀}

T5（Text-to-Text Transfer Transformer）将所有 NLP 任务统一建模为``文本到文本''的格式，并使用片段腐蚀作为预训练目标。

\begin{figure}[H]
\centering
\includegraphics[width=0.95\textwidth]{slides-images/slide-042.jpg}
\caption{片段腐蚀示例：输入 ``Thank you [X] me to your party [Y] week''，输出 ``[X] for inviting [Y] last'' ——模型需要生成被遮住的连续片段\protect\footnotemark}
\end{figure}
\footnotetext{Slides 第43页。}

片段腐蚀的工作方式：
\begin{enumerate}
    \item 在输入中随机遮住若干\textbf{连续片段}，用特殊标记（如 \texttt{[X]}、\texttt{[Y]}）替代
    \item Encoder 处理带遮罩的输入（保持双向上下文）
    \item Decoder 自回归\textbf{生成}被遮住的内容：先输出 \texttt{[X]}，再输出对应片段，再输出 \texttt{[Y]}，再输出对应片段...
\end{enumerate}

\begin{importantbox}{片段腐蚀 = BERT 的掩码思想 + 生成式输出}
片段腐蚀结合了两个世界的优点：
\begin{itemize}
    \item 像 BERT 一样利用\textbf{双向上下文}来理解输入
    \item 像语言模型一样以\textbf{自回归方式生成}输出
\end{itemize}
这使得预训练好的 T5 可以自然地应用于翻译、摘要等序列到序列任务。
\end{importantbox}

\subsection{T5 的``隐式知识检索''能力}

一个令人惊讶的发现是：T5 在预训练中看过大量类似``Franklin D. Roosevelt was born in [MASK]''的文本，之后经过在问答数据上微调，它能够\textbf{直接从参数中}回忆出答案，而非从外部知识库检索。

这种``闭卷''问答能力表明，预训练模型在其参数中\textbf{隐式存储了大量事实知识}。

\begin{warningbox}{模型回答总是看起来很流畅，但经常是错的}
Chris Manning 特别指出：这些模型的回答``always look very fluent, always look very reasonable, but they're frequently wrong.'' 这个问题（后来被称为``幻觉''，hallucination）至今仍是大语言模型的核心挑战之一。
\end{warningbox}

\subsection{本章小结}

Encoder-Decoder 预训练以 T5 为代表，通过片段腐蚀实现了双向理解与生成能力的统一。T5 将所有任务转化为``文本到文本''格式，提供了一种优雅的统一框架。模型在预训练过程中隐式存储了大量事实知识，但也伴随着幻觉问题。

%% ========== Part 5: Decoder 预训练（GPT 系列） ==========

